% =============================================================================
% event_nlos_survey.tex — Event Cameras for NLOS Communication: A Survey
% Companion to event_nlos_refs.bib (18 references)
% =============================================================================
\documentclass[10pt,a4paper]{article}

\usepackage[utf8]{inputenc}
\usepackage[T1]{fontenc}
\usepackage{cite}
\usepackage{amsmath,amssymb}
\usepackage{graphicx}
\usepackage{booktabs}
\usepackage{multirow}
\usepackage{url}
\usepackage{hyperref}
\usepackage{xcolor}
\usepackage{geometry}
\usepackage{array}
\usepackage{tabularx}

\geometry{margin=2.2cm}

\hypersetup{
  colorlinks=true,
  linkcolor=blue!60!black,
  citecolor=teal!70!black,
  urlcolor=blue!60!black
}

\newcommand{\stage}[1]{\textbf{Stage~#1}}
\newcommand{\etc}{\textit{etc.}\ }

\title{Event Cameras for Non-Line-of-Sight Optical Communication: \\
       A Survey of Technical Evolution and Future Directions}

\author{
  NLOS-OCC Research Group
}

\date{September 2026}

% =============================================================================
\begin{document}
\maketitle

% -----------------------------------------------------------------------------
\begin{abstract}
Event cameras (neuromorphic vision sensors, dynamic vision sensors) have
emerged as a transformative receiver technology for optical camera
communication (OCC), fundamentally addressing the temporal sampling
bottleneck that limits conventional frame-based cameras to sub-kHz
bandwidths. By asynchronously reporting pixel-level log-intensity changes
with microsecond resolution and a dynamic range exceeding 120\,dB, event
cameras enable communication links that are simultaneously high-rate,
motion-tolerant, and robust to extreme illumination---properties that are
especially valuable when the line-of-sight path is blocked and the
receiver must recover data from diffuse reflections. This survey traces
the technical evolution of event-camera-based optical communication from
its hardware origins (2008--2017) through four developmental stages:
proof-of-concept event-OCC (2017--2018), spatial encoding and beacon
systems (2022--2024), the first non-line-of-sight (NLOS) event-camera
demonstrations (2025), and long-range outdoor plus specialized
deployments (2025--2026). We organize the literature around three axes:
hardware capability, modulation/demodulation design, and channel
environment (line-of-sight vs.\ NLOS vs.\ underwater). For each stage we
distil the key technical mechanisms---OOK, N-pulse modulation, event
interval modulation, toggle demodulation, 8b/10b clock recovery, dynamic
QR-code spatial encoding, and adaptive reflection-based NLOS
schemes---and tabulate the experimentally demonstrated data-rate-vs-distance
frontier. We close with open challenges and a roadmap toward
event-camera-enabled NLOS communication networks for 6G-era applications.
\end{abstract}

\noindent\textbf{Index Terms}---Event camera, neuromorphic vision sensor,
dynamic vision sensor (DVS), optical camera communication (OCC),
non-line-of-sight (NLOS), visible light communication (VLC), rolling
shutter, modulation, clock recovery, 6G.

% =============================================================================
\section{Introduction}
\label{sec:intro}

Optical Camera Communication (OCC) uses an image sensor as the receiver
for visible light communication (VLC), exploiting the camera's spatial
resolution to discriminate among multiple LED transmitters and reject
ambient interference. The fundamental limitation of conventional
frame-based OCC is temporal: a commodity camera captures 30--60 frames
per second, so the effective sampling bandwidth is below 1\,kHz, even
with rolling-shutter stripe extraction. LED transmitters, by contrast,
are modulated in the MHz range. This sampling mismatch---three to four
orders of magnitude---caps the throughput of frame-based OCC and makes
it acutely sensitive to motion blur and frame-synchronization errors.

Event cameras, also called Dynamic Vision Sensors (DVS) or neuromorphic
vision sensors, dissolve this bottleneck. Each pixel operates
asynchronously and independently, generating an \emph{event}
$(x, y, t, p)$ only when the log-intensity change exceeds a threshold,
where $t$ is a microsecond-resolution timestamp and $p \in \{+1, -1\}$
is the polarity (brightening vs.\ dimming). The result is a sparse,
asynchronous stream with three properties that are ideal for OCC:
(i) \textbf{high temporal resolution} ($\sim$1\,$\mu$s, equivalent to
$\sim$20--1000\,kHz effective sampling), (ii) \textbf{wide dynamic range}
($>$120\,dB, vs.\ $\sim$60--70\,dB for CMOS cameras), and
(iii) \textbf{low data volume and latency} (only changing pixels are
transmitted, eliminating redundant background frames).

These properties become even more critical in Non-Line-of-Sight (NLOS)
OCC, where the direct LED-to-camera path is blocked and the receiver
must recover data from diffuse reflections off walls, the ground, or
ordinary objects. The reflected signal is drastically attenuated, smeared
by the bidirectional reflectance distribution function (BRDF) of the
surface, and superimposed on heterogeneous background illumination. An
event camera's high dynamic range can capture the faint reflected
intensity changes that a frame-based camera would bury in quantization
noise, and its asynchronous operation eliminates the frame-sync problem
that plagues reflected signals with unknown propagation delay.

\paragraph{Scope and contributions.}
This survey focuses specifically on event cameras as receivers for
optical communication, with emphasis on the progression toward NLOS
operation. We trace the technical evolution through four developmental
stages (Section~\ref{sec:evolution}), analyze the modulation and
demodulation design space (Section~\ref{sec:modulation}), examine the
NLOS event-camera channel in detail (Section~\ref{sec:nlos}), and close
with open challenges and a roadmap (Section~\ref{sec:roadmap}). We
summarize the experimentally demonstrated data-rate-vs-distance frontier
in Table~\ref{tab:frontier} and identify the key technical transitions
that enabled each advance.

% =============================================================================
\section{Hardware Foundations}
\label{sec:hardware}

The event-camera hardware evolution directly enabled the communication
performance gains surveyed in this paper. Three hardware milestones are
foundational:

\paragraph{DVS128 (2008).}
The first commercial event camera, introduced by Lichtsteiner
\emph{et al.}~\cite{lichtsteiner2008dvs}, featured a $128\times128$ pixel
array with 40\,$\mu$m pixel pitch, 120\,dB dynamic range, and 15\,$\mu$s
latency. Its $350$\,nm CMOS process produced high dark current, but the
asynchronous pixel architecture---a continuous-time logarithmic
photoreceptor coupled to a self-timed threshold comparator---established
the event-generation model that all subsequent sensors inherit. The
DVS128's modest resolution and high noise limited early event-OCC
experiments to short ranges and low data rates.

\paragraph{DAVIS240 (2014).}
Brandli \emph{et al.}~\cite{brandli2014davis240} introduced the DAVIS240,
a $240\times180$ sensor with 130\,dB dynamic range and 3\,$\mu$s latency.
Crucially, DAVIS240 interleaves event data with global-shutter
intensity frames, enabling hybrid processing pipelines that use frames
for transmitter localization and events for high-speed demodulation.

\paragraph{Samsung VGA DVS (2017).}
Son \emph{et al.}~\cite{son2017samsungdvs} demonstrated a $640\times480$
VGA-resolution DVS with 9\,$\mu$m pixels and 300\,Meps (million events
per second) throughput, fabricated in a BSI CIS process. The higher
spatial resolution and event throughput directly enabled longer-range
event-OCC links, as the receiver could resolve smaller LED images at
greater distances while maintaining sufficient event density for
demodulation. Subsequent sensors (e.g., Prophesee EVK4 with Sony IMX636,
$1280\times720$, 4.86\,$\mu$m pixel) continued this trend, reaching the
$\sim$1\,MHz per-pixel event-rate ceiling exploited by the latest
long-range demonstrations~\cite{nardo2026longrange}.

% =============================================================================
\section{Technical Evolution: Four Stages}
\label{sec:evolution}

The application of event cameras to optical communication has progressed
through four identifiable stages, each defined by a key technical
transition. Table~\ref{tab:frontier} summarizes the experimentally
demonstrated performance milestones.

\begin{table*}[t]
\centering
\caption{Experimentally demonstrated event-OCC performance frontier.}
\label{tab:frontier}
\small
\begin{tabularx}{\textwidth}{l l l l l X}
\toprule
\textbf{Year} & \textbf{Stage} & \textbf{Data Rate} & \textbf{Distance} & \textbf{BER} & \textbf{Key Technique} \\
\midrule
2017 & \stage{1} & 500 bit/s & 5 m & $<10^{-4}$ & OOK + DVS128 receiver~\cite{perez2017eventocc} \\
2018 & \stage{1} & 16 kbit/s & 10 m & N/A & i-PPM + FSK multi-level~\cite{shen2018eventocc} \\
2022 & \stage{2} & 500 bit/s & 100 m & 0 (error-free) & Smart visual beacon + LED~\cite{wang2022beacon} \\
2024 & \stage{2} & 2.5 kbit/s & 0.5 m & 94\% acc. & Dynamic environment framework~\cite{vonarnim2024dynamic} \\
2024 & \stage{2} & --- & --- & --- & N-pulse modulation~\cite{aranda2024npulse} \\
\midrule
2025 & \stage{3} & 114 kbit/s & 1.25 m & 1\% & Dynamic QR spatial encoding~\cite{su2025motionaware} \\
2025 & \stage{3} & --- & indoor & --- & First NLOS event-camera OCC~\cite{nishar2025nlosneuromorphic} \\
\midrule
2025 & \stage{4} & 60 kbit/s & 200 m & $<10^{-3}$ & OOK + toggle + digital PLL~\cite{sumino2025longrange} \\
2025 & \stage{4} & 30 kbit/s & 400 m & $<10^{-3}$ & (outdoor, same system)~\cite{sumino2025longrange} \\
2025 & \stage{4} & 28 kbit/s & 10 m & --- & Event Interval Modulation (EIM)~\cite{sumino2025eim} \\
2026 & \stage{4} & 80 kbit/s & 320 m & $<10^{-3}$ & 8b/10b + OOK-NRZ + PLL, 120$^\circ$ FoV~\cite{nardo2026longrange} \\
2026 & \stage{4} & 240 kbit/s & 320 m & $<10^{-3}$ & Multi-node (3 beams, same system)~\cite{nardo2026longrange} \\
\bottomrule
\end{tabularx}
\end{table*}

% -----------------------------------------------------------------------------
\subsection{\stage{1}: Proof-of-Concept Event-OCC (2017--2018)}
\label{sec:stage1}

The first demonstrations that an event camera could serve as an OCC
receiver appeared in 2017--2018. Perez-Ramirez \emph{et al.}
\cite{perez2017eventocc} demonstrated a 500\,bit/s link over 5\,m with
BER $<10^{-4}$ using On-Off Keying (OOK) and a DVS128 receiver. The key
insight was that an OOK-modulated LED produces a rising edge (ON event)
on the pixels imaging the LED when it turns on and a falling edge (OFF
event) when it turns off; the event stream thus directly encodes the
bit sequence without the rolling-shutter stripe extraction required by
frame-based cameras.

Shen \emph{et al.}~\cite{shen2018eventocc} extended this to multi-level
modulation, achieving 16\,kbit/s at 10\,m using inverse Pulse Position
Modulation (i-PPM) combined with Frequency Shift Keying (FSK). This work
first exploited the event camera's microsecond temporal resolution for
frequency-domain encoding, moving beyond binary OOK. However, BER was
not evaluated, and the system required a static, aligned link---no motion
or NLOS capability was demonstrated.

\paragraph{Stage 1 summary.}
These works established that event cameras can demodulate optical
signals at rates far exceeding the frame-rate limit, but the
demonstrated throughput ($\leq$16\,kbit/s) and range ($\leq$10\,m)
remained modest, and all links were line-of-sight and static.

% -----------------------------------------------------------------------------
\subsection{\stage{2}: Spatial Encoding and Beacon Systems (2022--2024)}
\label{sec:stage2}

The second stage introduced spatial encoding---using the event camera's
2D pixel array as a communication resource, not just a single-pixel
detector---and extended the link range by an order of magnitude.

Wang \emph{et al.}~\cite{wang2022beacon} proposed a ``smart visual
beacon'' architecture at IROS 2022, combining an LED transmitter with
event-camera demodulation algorithms. They achieved 100\,m at 500\,bit/s
with zero reported bit errors, quantitatively evaluating the tradeoff
among LED transmission rate, communication distance, and message
accuracy. The beacon architecture was designed for robotic
applications---the event camera on a robot demodulates a beacon on a
stationary fixture---and it demonstrated that event-OCC could operate at
ranges relevant to indoor navigation and V2X scenarios.

Aranda \emph{et al.}~\cite{aranda2024npulse} introduced N-pulse
modulation, a scheme specifically designed for the event camera's
operating principles. Rather than encoding each bit as a single intensity
transition, N-pulse encodes a symbol as a group of $N$ pulses with a
defined pattern, adapting the encoding to the packet's bit composition.
This significantly reduced demodulation complexity compared to OOK and
PPM, as the receiver could exploit the temporal structure of the pulse
group rather than relying on per-edge thresholding. The same group
subsequently evaluated pulse-shaped waveform modulation
\cite{aranda2025pulseshaped}, leveraging the sensor's high temporal
resolution for shaped pulses that improve noise resilience.

Von Arnim \emph{et al.}~\cite{vonarnim2024dynamic} addressed the
practical challenge of dynamic environments, where ambient light and
target motion change during communication. Their framework achieved
2.5\,kbit/s with 94\% message accuracy at 0.5\,m---modest numbers, but
the contribution was robustness to changing conditions rather than raw
performance.

\paragraph{Stage 2 summary.}
Stage 2 expanded event-OCC from a single-pixel novelty to a spatially
encoded, robot-relevant communication modality. Range extended to 100\,m,
but throughput remained below 10\,kbit/s, and all links were still
line-of-sight.

% -----------------------------------------------------------------------------
\subsection{\stage{3}: NLOS Event-Camera Communication (2025)}
\label{sec:stage3}

The third stage---the central focus of this survey---marks the first
demonstrations of event cameras for NLOS communication, where the signal
is recovered from reflections rather than direct line-of-sight.

\subsubsection{First NLOS event-camera OCC}
Nishar \emph{et al.}~\cite{nishar2025nlosneuromorphic} (EWSN 2025,
arXiv:2503.11226) demonstrated the first NLOS optical wireless
communication using a neuromorphic camera. The key innovation was
leveraging \emph{passive reflections from everyday objects} as the
communication channel: an LED modulated with OOK or N-pulse modulation
illuminates a surface (e.g., a wall, a poster, furniture), and the event
camera detects the intensity changes reflected from that surface---even
though the camera has no direct view of the LED. The system overcomes
the event camera's fundamental limitation (it requires intensity
\emph{changes} to generate events; a static scene produces no events) by
deliberately inducing VLC-modulated intensity variations, thereby
generating events even under zero-motion conditions.

The paper introduced an \emph{adaptive N-pulse modulation} scheme that
dynamically adjusts the encoding based on the packet's bit composition,
optimizing both data rate and robustness. Experimental results revealed
that lighter-colored, glossy objects are better NLOS communication
surfaces (higher reflectivity, more specular reflection), while larger
objects and matte finishes produce higher error rates due to multipath
reflections that smear the temporal structure of the events.

\subsubsection{Motion-aware spatial encoding}
Su \emph{et al.}~\cite{su2025motionaware} (IEEE RAL 2025,
arXiv:2412.00816) introduced a complementary advance: a dynamic visual
marker (effectively a QR-code-like spatial pattern displayed on a
screen) combined with event-based tracking and decoding algorithms. The
event camera tracks the marker's corners and decodes the spatially
encoded data simultaneously, achieving 114\,kbit/s in static
conditions---the highest short-range event-OCC throughput to date---and
1\,cm localization accuracy with 1\% BER under various camera motions.
The critical contribution for NLOS is that the event camera's
motion-blur immunity enables communication even when the receiver is
moving, a scenario that defeats frame-based OCC due to blur-induced
stripe distortion. The open-source code and dataset are available at
\url{https://github.com/suhang99/EventOCC}.

\paragraph{Stage 3 summary.}
Stage 3 established that event cameras can communicate via NLOS
reflections and under receiver motion---the two capabilities that
frame-based cameras cannot provide. Throughput reached 114\,kbit/s
(short-range, spatial encoding), but NLOS-through-reflection throughput
remained low (order of kbit/s), limited by the reflection channel's
attenuation and multipath.

% -----------------------------------------------------------------------------
\subsection{\stage{4}: Long-Range and Specialized Deployments (2025--2026)}
\label{sec:stage4}

The fourth stage, currently unfolding, pushes event-OCC to
deployment-relevant ranges and environments: outdoor long-range links,
underwater communication, and multi-node networks.

\subsubsection{Long-range outdoor}
Sumino \emph{et al.}~\cite{sumino2025longrange} (arXiv:2510.14266,
GLOBECOM 2025) achieved the first sub-$10^{-3}$ BER long-range outdoor
event-OCC: 60\,kbit/s at 200\,m and 30\,kbit/s at 400\,m, using OOK with
toggle demodulation and a digital phase-locked loop (PLL) for clock
recovery. The toggle demodulation scheme exploits the event camera's
polarity output: each OOK symbol produces an ON-event burst followed by
an OFF-event burst, and the demodulator toggles its bit decision based
on the event polarity, providing intrinsic noise rejection. The digital
PLL tracks the symbol clock from the event timestamps, handling the
timing jitter that outdoor turbulence introduces.

Nardo \emph{et al.}~\cite{nardo2026longrange} (arXiv:2609.29776) pushed
this further: 80\,kbit/s at 320\,m with a 120$^\circ$ field-of-view
fisheye lens, using 8b/10b line coding with OOK-NRZ and a Bang-Bang PLL.
The 8b/10b coding guarantees at least one edge transition every five
symbols, ensuring continuous event generation even for long runs of
identical bits. In multi-node operation, three parallel beams yielded
240\,kbit/s aggregate, with a real-time processing latency under 10\,ms
per 50\,ms event batch. The system used a Prophesee EVK4 (Sony IMX636)
sensor; the authors identified the sensor's readout bottleneck
($\sim$100 lines/$\mu$s) as the factor preventing the theoretical
1\,MHz bandwidth from being reached in practice.

\subsubsection{Novel modulation: Event Interval Modulation}
Sumino \emph{et al.}~\cite{sumino2025eim} (arXiv:2510.14245) proposed
Event Interval Modulation (EIM), the first modulation scheme designed
specifically for event cameras. Rather than encoding information in the
presence/absence of events (OOK) or in pulse patterns (N-pulse), EIM
encodes information in the \emph{time intervals between consecutive
events}. This fully exploits the event camera's asynchronous temporal
resolution and enables higher-order modulation: the interval length
maps to a symbol, so $M$-ary EIM carries $\log_2 M$ bits per event.
The proof-of-concept achieved 28\,kbit/s at 10\,m and 8.4\,kbit/s at
50\,m indoors, setting a new benchmark for event-OCC bit rate in the
moderate-range regime.

\subsubsection{Underwater event-camera ISAC}
Fang \emph{et al.}~\cite{fang2026underwater} (Frontiers in Marine
Science 2026) extended event-camera OCC to the underwater domain,
demonstrating an event-based camera-enabled underwater optical
integrated sensing and communication (ISAC) system. This is notable as
the first cross-medium event-OCC: the underwater channel's
wavelength-dependent absorption and particulate scattering are
analogous to an extreme NLOS channel, and the event camera's high
dynamic range is essential for coping with the rapid intensity
fluctuations caused by water turbulence.

\subsubsection{MISO and real-time systems}
Mu\~noz \emph{et al.}~\cite{munoz2026miso} (IEEE PTL 2026) designed a
feedback-aware MISO event-OCC system with an LED array transmitter and
event camera receiver, achieving near-real-time operation by leveraging
event-driven sensing and a lightweight demodulation scheme. The feedback
channel allows the transmitter to adapt its modulation parameters based
on receiver-side event-rate feedback, closing the loop for adaptive
event-OCC.

\paragraph{Stage 4 summary.}
Stage 4 has brought event-OCC to deployment-relevant regimes:
320\,m outdoor range, 240\,kbit/s multi-node throughput, underwater
operation, and closed-loop adaptation. The data-rate-vs-distance
frontier has advanced by 2--3 orders of magnitude since Stage 1.

% =============================================================================
\section{Modulation and Demodulation Design Space}
\label{sec:modulation}

The event camera's unique sensing model---asynchronous, differential,
polarity-labeled---has motivated a distinct modulation design space that
differs fundamentally from both photodiode-based VLC and frame-based OCC.
We organize the schemes into three categories.

\subsection{Edge-Transition-Based Schemes}
\label{sec:edge}

The simplest and most widely used approach exploits the event camera's
native output: each intensity transition generates an event. OOK
modulation maps a ``1'' to an ON-event and a ``0'' to an OFF-event; the
receiver demodulates by checking the polarity of the event stream within
each symbol period~\cite{perez2017eventocc, sumino2025longrange,
nardo2026longrange}. The challenge is clock recovery: long runs of
identical bits produce no events, causing the receiver to lose symbol
alignment. Two solutions have been demonstrated:
\begin{itemize}
  \item \textbf{8b/10b line coding}~\cite{nardo2026longrange}: guarantees
    at least one transition every five symbols, ensuring continuous
    event generation at the cost of 25\% overhead.
  \item \textbf{Toggle demodulation}~\cite{sumino2025longrange}: each
    symbol period is expected to contain both an ON and an OFF event;
    the demodulator toggles its decision based on which arrives first,
    providing intrinsic noise rejection.
\end{itemize}

\subsection{Pulse-Pattern Schemes}
\label{sec:pulse}

N-pulse modulation~\cite{aranda2024npulse} encodes a symbol as a group
of $N$ pulses with a defined temporal pattern, adapting the pattern to
the bit composition of the packet. The receiver correlates the event
stream against the expected patterns, exploiting the temporal structure
for noise rejection. The adaptive variant~\cite{nishar2025nlosneuromorphic}
dynamically selects the pattern based on the packet's bit distribution,
optimizing the tradeoff between data rate and robustness. Pulse-shaped
waveform modulation~\cite{aranda2025pulseshaped} further shapes the
pulse waveform to exploit the sensor's high temporal resolution.

\subsection{Interval-and-Position Schemes}
\label{sec:interval}

Event Interval Modulation (EIM)~\cite{sumino2025eim} encodes information
in the time intervals between consecutive events, enabling $M$-ary
modulation where each interval maps to $\log_2 M$ bits. Inverse PPM
\cite{shen2018eventocc} and FSK~\cite{shen2018eventocc} encode
information in the temporal position and frequency of event bursts,
respectively. These schemes fully exploit the microsecond temporal
resolution but require precise timing extraction and are sensitive to
event jitter.

\subsection{Spatial Encoding}
\label{sec:spatial}

The event camera's 2D pixel array enables spatial multiplexing: the
dynamic QR-code approach~\cite{su2025motionaware} displays a
spatially-encoded pattern on a screen, and the event camera decodes the
pattern by tracking corner features and reading the spatial cells. This
achieves the highest short-range throughput (114\,kbit/s) but requires a
display transmitter and is limited by the camera's spatial resolution.
The smart visual beacon~\cite{wang2022beacon} uses a simpler spatial
structure (a single LED or small array) with spatial tracking for
robotic applications.

% =============================================================================
\section{The NLOS Event-Camera Channel}
\label{sec:nlos}

NLOS event-camera communication is the frontier that most distinguishes
this technology from photodiode-based VLC and frame-based OCC. We
examine the channel characteristics and the techniques that make NLOS
event-OCC feasible.

\subsection{The Reflection Channel}
\label{sec:reflection}

In NLOS-OCC, the LED transmitter illuminates a reflecting surface (wall,
ground, object), and the event camera images the surface to detect the
modulated intensity changes in the reflected light. The channel is
characterized by:
\begin{itemize}
  \item \textbf{Path loss}: the reflected optical power is attenuated by
    the product of the LED-to-surface and surface-to-camera path losses,
    plus the surface reflectance (typically 0.05--0.9 depending on
    material and wavelength).
  \item \textbf{BRDF smearing}: the Bidirectional Reflectance
    Distribution Function of the surface spreads the reflected light
    over a range of angles, reducing the intensity at the camera and
    blurring the temporal modulation. Matte surfaces produce diffuse
    (Lambertian) reflections with severe multipath; glossy surfaces
    produce more specular reflections that preserve temporal structure.
  \item \textbf{Multipath}: photons reflected via different paths arrive
    at different times, causing inter-symbol interference. Nishar
    \emph{et al.}~\cite{nishar2025nlosneuromorphic} confirmed that
    larger objects and matte finishes produce higher BER due to
    multipath, while lighter-colored glossy surfaces are better.
\end{itemize}

A simplified mathematical model for the NLOS-OCC channel with a patterned
reflective surface was developed by He \emph{et al.}~\cite{he2023nlosmodel},
providing the channel-characterization foundation for NLOS event-camera
work.

\subsection{Why Event Cameras Excel at NLOS}
\label{sec:why}

Three properties of event cameras are particularly advantageous for NLOS
communication:
\begin{enumerate}
  \item \textbf{Dynamic range} ($>$120\,dB): the reflected signal may be
    30--60\,dB weaker than the direct signal. A frame-based camera with
    60--70\,dB dynamic range cannot simultaneously capture the bright
    background and the faint reflected modulation without saturation or
    quantization loss. An event camera's logarithmic photoreceptor
    handles this range natively.
  \item \textbf{Asynchronous operation}: the reflected signal's
    propagation delay is unknown and varies with geometry. Frame-based
    OCC requires frame-level synchronization, which fails when the
    delay is not an integer multiple of the frame period. Event cameras
    have no frame structure; each event is timestamped independently,
    so the receiver recovers timing from the event stream itself via a
    digital PLL~\cite{sumino2025longrange, nardo2026longrange}.
  \item \textbf{Motion tolerance}: in mobile NLOS scenarios (e.g., a
    vehicle receiving a reflected signal from a taillight while
    moving), frame-based cameras suffer motion blur that smears the
    stripes. Event cameras produce no blur---each pixel reports the
    intensity change at the moment it occurs, regardless of camera
    motion~\cite{su2025motionaware}.
\end{enumerate}

\subsection{The Static-Scene Challenge}
\label{sec:static}

A unique challenge of event cameras is that a static scene produces no
events. In LOS OCC this is not a problem (the LED is modulated, so its
pixels always have intensity changes), but in NLOS OCC the reflected
intensity change may be too small to exceed the event threshold,
especially with matte surfaces. Nishar \emph{et al.}
\cite{nishar2025nlosneuromorphic} addressed this by deliberately
inducing VLC-modulated intensity variations strong enough to generate
events via reflection, and by tuning the event threshold to the
reflection's contrast level. Future work may exploit adaptive threshold
control---dynamically adjusting the event threshold per-pixel based on
the local reflected signal strength.

% =============================================================================
\section{Open Challenges and Roadmap}
\label{sec:roadmap}

\paragraph{NLOS throughput.}
The current NLOS event-OCC throughput (order of kbit/s via passive
reflection) is 1--2 orders of magnitude below the LOS frontier
(80--240\,kbit/s). Closing this gap requires: (i) higher-gain optical
front-ends (telescopic lenses, Fresnel concentrators) to collect more
reflected photons; (ii) adaptive event-threshold control to match the
reflection's contrast; and (iii) channel coding (Raptor/RaptorQ fountain
codes) tailored to the burst-error profile of the reflection channel.

\paragraph{NLOS channel modeling.}
Accurate analytical and ray-tracing models for the event-camera NLOS
channel---incorporating BRDF, multipath, and the event-generation
threshold---are lacking. The simplified model of He \emph{et al.}
\cite{he2023nlosmodel} is a start but does not account for the event
camera's differential sensing model. A unified channel model that
connects the optical power at the pixel to the event-generation rate is
needed.

\paragraph{Modulation for NLOS.}
The adaptive N-pulse scheme~\cite{nishar2025nlosneuromorphic} is the
only modulation designed for NLOS event-OCC. EIM~\cite{sumino2025eim}
and 8b/10b+OOK~\cite{nardo2026longrange} have been demonstrated only in
LOS. Extending these to NLOS---where the signal-to-noise ratio is much
lower and multipath causes inter-event interference---is an open problem.

\paragraph{Multi-node and networking.}
The multi-node demonstration of Nardo \emph{et al.}
\cite{nardo2026longrange} (3 beams, 240\,kbit/s) is the only multi-node
event-OCC result. A complete networking stack---medium access control,
addressing, relaying---for event-camera networks does not yet exist.
The IEEE 802.15.7a-2024 standard's AI-PHY and relaying provisions
provide a framework, but event-camera-specific adaptations are needed.

\paragraph{Cross-medium and extreme environments.}
The underwater event-camera ISAC of Fang \emph{et al.}
\cite{fang2026underwater} is a promising start, but the underwater
channel's absorption and turbulence are far more severe than air.
Space-air-ground-sea architectures integrating event-camera OCC with RF
and FSO are a long-term vision.

\paragraph{Hardware co-design.}
The readout bottleneck identified by Nardo \emph{et al.}
\cite{nardo2026longrange} ($\sim$100 lines/$\mu$s) limits the practical
bandwidth to well below the sensor's per-pixel event-rate ceiling.
Sensor-level innovations---parallel readout, on-chip event
processing---are needed to unlock the full $\sim$1\,MHz bandwidth.

% =============================================================================
\section{Conclusion}
\label{sec:conclusion}

Event cameras have progressed from a vision-sensing novelty to a
communication receiver capable of 80\,kbit/s at 320\,m (LOS outdoor),
114\,kbit/s at short range (spatial encoding), and---crucially---kbit/s-rate
communication via passive NLOS reflections from everyday objects. The
technical evolution has advanced through four stages: proof-of-concept
OOK (2017), spatial encoding and beacon systems (2022--2024), the first
NLOS event-camera demonstrations (2025), and long-range outdoor plus
specialized deployments (2025--2026). The key enablers have been the
event camera's microsecond temporal resolution, $>$120\,dB dynamic range,
and asynchronous, motion-blur-free operation---properties that are
uniquely suited to the NLOS reflection channel. The open challenges are
NLOS throughput, channel modeling, NLOS-specific modulation, multi-node
networking, and hardware co-design. As event-camera hardware continues
to improve and the 6G vision of integrated sensing and communication
matures, event-camera-based NLOS OCC is positioned to become a key
enabling technology for vehicular, robotic, and IoT networks in
cluttered, occluded, and extreme environments.

% =============================================================================
\bibliographystyle{plain}
\bibliography{event_nlos_refs}

\end{document}
